\lecture{4}{Thu 15 Jun 2023 05:08}{General Order Moment Tail Bound}

\subsection{Moment Bound of General Order}

In this section we try to generalize the moment bound in the previous section (variance bound) to higher order. Surprisingly all arguments generalize to higher order with little additional difficulty. 

The result is that for all $\delta>0$ small,

\begin{equation*}
	P \left( \left| \sum_y \Delta_y^{(x,m)} - \rho_y^{(x,m)} \right| > \varepsilon \sqrt{n} \right) 
	= O\left( n^{\frac{1}{2}-\frac{1}{2}K p + \delta } \right)
.\end{equation*}


\subsubsection{Implication of the bound}

The bound improves as $K$ increases; and by taking large enough $K$ we can get a bound better than any (general) polynomial in $n$.

\subsubsection{Derivation of the bound}

\textbf{Step 1. Moment of sectional drift}

By ``sectional drift'' we mean the drift attained at $y$ between two consecutive left excursions.

\begin{lemma}
	(Drift at single step)
	\[
		\mathbb{E}\left[ \left| \alpha(i, \tau_i^{\mathcal{B}} - i) \right| ^K \right] 
		= O\left(i^{-Kp - \frac{1}{2}K}\right)
	.\] 
\end{lemma}
\begin{proof}
	Minor adaptation of Lemma~\ref{lem:sectional-drift-1} and the one following.
\end{proof}

\begin{lemma}
	(Control of sectional drift)
	\[
		\mathbb{E}\left[ \left|   
			\sum_{j = \tau_{i-1}^{\mathcal{B}}-i+1}^{\tau_i^{\mathcal{B}}-i}   
			\alpha(i, j) 
		\right| ^K \right]
		= O\left(i^{-Kp - \frac{1}{2}K}\right)
	.\] 
\end{lemma}
\begin{proof}
	Intuition: $\tau_i - \tau_{i-1} = O(1)$ as given by Lemma~\ref{lem:section-concentration}. Also, $O(1)$ perturbation in $j$ results in lower order change in $\alpha(i,j)$.
\end{proof}


\textbf{Step 2. Moment of local drift}

\begin{lemma}
	\[
		\mathbb{E}\left[ \left| \Delta_y^{(x,m)}\right|^K  \right] =  
		\mathbb{E}\left[ \left| 
		\sum_{i = 1}^{\mathcal{E}_y^{(x,m)}} 
			\sum_{j = \tau_{i-1}^{\mathcal{B}}-i+1}^{\tau_i^{\mathcal{B}}-i}   
			\alpha(i, j) 
	\right| ^K \right] 
		= O\left(\mathcal{E}^{-Kp + \frac{1}{2}K}\right)
	.\]
	\label{lem:local-drift-moment}
\end{lemma}
\begin{proof}
	For simplicity of notation, write $\alpha_i$ for the inner sum.     

	\begin{align*}
		\mathbb{E}\left[ \left| \sum_{i = 1}^{\mathcal{E}} \alpha_i \right| ^K \right] 
		&= \mathbb{E}\left[ 
			\left|
			\sum_{\left( i_m \right)_{m = 1}^K \in \{1, \ldots, \mathcal{E}\}^K} 
			\prod_{m = 1}^{K} \alpha_{i_m} 
		\right|
		\right] \\
		&\le  
		\sum_{\left( i_m \right)_{m = 1}^K \in \{1, \ldots, \mathcal{E}\}^K} 
		\mathbb{E}\left[ \left| \prod_{m = 1}^{K} \alpha_{i_m} \right| \right]
	.\end{align*}

	By generalized Holder Inequality and the previous lemma, for any $\left( i_m \right) _{m = 1}^K$,
	\begin{align*}
		\mathbb{E}\left[ \left| \prod_{m = 1}^{K} \alpha_{i_m} \right| \right]
		&\le \prod_{m = 1}^{K} \left(
			\mathbb{E}\left[
			|\alpha_{i_m}|^K
		\right] \right)^{\frac{1}{K}}  \\
		&= O(i^{-Kp - \frac{1}{2}K}) 
	.\end{align*}
	where, $i$ stands for the minimum of all the $i_m$.
	Finally, considering all $K$-tuples, we obtain
\[
		\mathbb{E}\left[ \left| \sum_{i = 1}^{\mathcal{E}} \alpha_i \right| ^K \right] 
		\le 
		\sum_{\left( i_m \right)_{m = 1}^K \in \{1, \ldots, \mathcal{E}\}^K}  
		C \min(i_m)^{- K p - \frac{1}{2} K}
		= O\left(\mathcal{E}^{-Kp + \frac{1}{2}K}\right)
.\] 
\end{proof}
\begin{note}{Note}
	The last equality may require a more careful analysis.
\end{note}




\textbf{Step 3. Apply $O(\sqrt{n}) $ estimate of $\mathcal{E}$}

We may let $\mathcal{E} = O(\sqrt{n}) $ and $\sum_y 1  = O(\sqrt{n} )$, though these need to be carefully justified later. Hence 


\[
	\mathbb{E}\left[ \left| \Delta_y^{(x,m)}\right|^K  \right]
		= O\left(n^{-\frac{1}{2}Kp + \frac{1}{4}K}\right)
.\] 

\textbf{Step 4. Control the martingale components}

\begin{note}{Note}
	We need control of $\mathbb{E}\left[  \left| \Delta_y^{(x,m)} - \rho_y^{(x,m)} \right| ^K\right] $ in terms of $\mathbb{E}\left[ \left| \Delta_y^{(x,m)} \right| ^K \right] $.
	
\end{note}
\[
	\mathbb{E}\left[ \left| \Delta_y^{(x,m)} - \rho_y^{(x,m)}\right|^K  \right]
		= O\left(n^{-\frac{1}{2}Kp + \frac{1}{4}K}\right)
.\]

\textbf{Step 5. Bound the martingale sum}

From control of martingale components and moment tail bound (Proposition~\ref{prop:moment-bound}), we obtain for specific $y, x$,

\[
	P\left( \left| \Delta_y^{(x,m)} - \rho_y^{(x,m)}\right| > \varepsilon n^{\frac{1}{4} - \frac{\kappa}{2}}   \right) \le O\left( n^{-\frac{1}{2} K p + \frac{1}{2} K \kappa} \right) 
.\] 



When $p > \frac{1}{K}$ and take $\kappa \in (0,p - \frac{1}{k})$, we can apply relaxed Azuma Inequality (Azuma Inequality applied to a good event where the martingale is Lipchitz) and obtain a bound similar to $\ref{eqn:azuma-bound-1}$.

\[
	\sum_y P\left( \left| \Delta_y^{(x,m)} - \rho_y^{(x,m)}\right| > \varepsilon n^{\frac{1}{4} - \frac{\kappa}{2}}   \right) 
	\le O\left( n^{\frac{1}{2}-\frac{1}{2} K p + \frac{1}{2} K \kappa} \right) 
.\]

 \begin{multline}
	P \left( \left| \sum_y \Delta_y^{(x,m)} - \rho_y^{(x,m)} \right| > \varepsilon \sqrt{n} \right) \\
	\le O\left( n^{\frac{1}{2}-\frac{1}{2}K p + \frac{1}{2} K \kappa } \right) + O\left( \exp\left( -\frac{ n^{\kappa}}{2}  \right)  \right) 
	= O\left( n^{\frac{1}{2}-\frac{1}{2}K p + \frac{1}{2} K \kappa } \right)
.\end{multline} 



\newpage

\subsection{Appendix}

\begin{proposition}
	(general order moment bound)
	\[
		P(|X| > a) \le \frac{\mathbb{E}\left[ |X|^K \right] }{a^K}
	.\] 
	\label{prop:moment-bound}
\end{proposition}
\begin{proof}
	Apply Chebyshev inequality with $\varphi(x) = x^K$.
\end{proof}


\begin{theorem}[Generalized Holder's Inequality]
	(\cite{finnerHolderInequality1992})
	Let $(\Omega, \mathscr{A}, \mu)$ be a measure space and let $L_p(\Omega, \mathscr{A}, \mu)$ be the set of $p$-integrable $(1 \leq p<\infty)$ measurable functions from $(\Omega, \mathscr{A})$ into $(\mathbb{R}, \mathscr{B})$, where $\mathscr{B}$ denotes the Borel $\sigma$-field over $\mathbb{R}$. Let $m \geq 2, p_j>1$ with $\sum_{j=1}^m p_j^{-1}=1$ and let $f_j \in L_{p_j}(\Omega, \mathscr{A}, \mu), j=1, \ldots, m$. Then it is well known [see, e.g., Kufner, John and Fučik (1977), page 67] that $\Pi_{j=1}^m\left|f_j\right| \in L_1(\Omega, \mathscr{A}, \mu)$ and
\[
\int \prod_{j=1}^m\left|f_j\right| d \mu \leq \prod_{j=1}^m\left(\int\left|f_j\right|^{p_j} d \mu\right)^{1 / p_j}
\]
with equality iff either there exists at least one $j$ with $\int\left|f_j\right|^{p_j} d \mu=0$ or there exist constants $A_j>0, j=1, \ldots, m$, such that
\[
A_1\left|f_1\right|^{p_1}=\cdots=A_m\left|f_m\right|^{p_m}[\mu] .
\]
\label{thm:general-holder}
\end{theorem}


